\section{Experimental Setup}
\label{sec:experimental_setup}

\subsection{Benchmark Functions}
\label{subsec:benchmarks}

The experimental evaluation uses six classical benchmark functions representing two structurally distinct groups: high-dimensional problems with regular landscapes and low-dimensional problems with rugged, deceptive landscapes.

The first group includes three functions evaluated in $d = 100$ dimensions. The Sphere function~(F1) serves as a unimodal baseline with a smooth, convex landscape. Rosenbrock's function~(F5) provides a unimodal challenge characterized by a narrow, curved valley where convergence is highly sensitive to search direction. Griewank's function~(F11) introduces regular multimodality through oscillatory components superimposed on a quadratic base.

The second group comprises three variants of the Shekel function in a four-dimensional space ($d = 4$), differing in the number of components~$m$: F21~($m{=}5$), F22~($m{=}7$), and F23~($m{=}10$). The Shekel family is characterized by closely spaced deep optima, making it particularly challenging for population-based methods prone to premature convergence.

All functions were implemented using the \texttt{Opfunu} library~\cite{Van_Thieu_2024_Opfunu}, ensuring reproducibility. Table~\ref{tab:benchmark_functions} summarizes the selected benchmarks.

\begin{table*}[!ht]
    \caption{Benchmark functions used in the experimental evaluation.}
    \label{tab:benchmark_functions}
    \centering
    \begin{tabular}{lcccc}
    \toprule
    Function & Type & Dim. & Bounds & Global Opt. \\
    \midrule
    F1 (Sphere)            & Unimodal   & $100$ & $[-100, 100]^d$ & $0.0$ \\
    F5 (Rosenbrock)        & Unimodal   & $100$ & $[-30, 30]^d$   & $0.0$ \\
    F11 (Griewank)         & Multimodal & $100$ & $[-600, 600]^d$ & $0.0$ \\
    F21 (Shekel, $m{=}5$)  & Multimodal & $4$   & $[0, 10]^d$     & $-10.1532$ \\
    F22 (Shekel, $m{=}7$)  & Multimodal & $4$   & $[0, 10]^d$     & $-10.4029$ \\
    F23 (Shekel, $m{=}10$) & Multimodal & $4$   & $[0, 10]^d$     & $-10.5363$ \\
    \bottomrule
    \end{tabular}
\end{table*}

\subsection{Parameter Settings and Evaluation Protocol}
\label{subsec:protocol}

All algorithmic variants were evaluated under identical conditions. The population size was fixed at $50$~particles and the maximum number of iterations was set to $500$. No additional stopping criteria were applied.

Each algorithm--function combination was executed over $31$ independent runs. A shared seeding strategy was adopted: all algorithms used the same random seed for the $k$-th run, with the initial seed set to $42$ and incremented by one for each subsequent run, yielding seeds $\{42, 43, \ldots, 72\}$.

The baseline PSO with linearly decreasing inertia weight was evaluated alongside eight fuzzy-controlled variants (PSO--FCS). These variants result from the full factorial combination of the four input set configurations described in Section~\ref{subsec:input_mf_configs} (I1--I4) and two granularity levels (3 and 5~labels), with the output fuzzy set~(A) and rule base held constant. Table~\ref{tab:algorithm_configs} summarizes the nine configurations, totaling $9 \times 6 \times 31 = 1{,}674$ independent experiments.

\begin{table*}[!ht]
    \caption{Algorithm configurations evaluated. All PSO--FCS variants share the same rule base and output set; only input membership functions and granularity vary.}
    \label{tab:algorithm_configs}
    \centering
    \begin{tabular}{lcccc}
    \toprule
    Algorithm & $w$ Control & Output Set & Labels & Input Set \\
    \midrule
    PSO (baseline)   & Linear decrease & --- & --- & --- \\
    PSO--FCS--I1--3L & Fuzzy (FCS) & A & 3 & I1 (Standard) \\
    PSO--FCS--I2--3L & Fuzzy (FCS) & A & 3 & I2 (Narrow) \\
    PSO--FCS--I3--3L & Fuzzy (FCS) & A & 3 & I3 (Wide) \\
    PSO--FCS--I4--3L & Fuzzy (FCS) & A & 3 & I4 (Shoulder) \\
    PSO--FCS--I1--5L & Fuzzy (FCS) & A & 5 & I1 (Standard) \\
    PSO--FCS--I2--5L & Fuzzy (FCS) & A & 5 & I2 (Narrow) \\
    PSO--FCS--I3--5L & Fuzzy (FCS) & A & 5 & I3 (Wide) \\
    PSO--FCS--I4--5L & Fuzzy (FCS) & A & 5 & I4 (Shoulder) \\
    \bottomrule
    \end{tabular}
\end{table*}

Performance was assessed using the final best fitness value and the relative gap to the known global optimum, defined as
\begin{equation}
    \text{Gap}(\%) =
    \frac{\left| f_{\text{best}} - f_{\text{opt}} \right|}
    {\left| f_{\text{opt}} \right|} \times 100.
    \label{eq:gap_metric}
\end{equation}

Robustness was evaluated via the mean and standard deviation of final fitness values across the 31~runs. Statistical significance was assessed using the Wilcoxon rank-sum test at $\alpha = 0.05$, complemented by the Mann--Whitney~U test as a non-paired confirmation. Computational cost was measured using wall-clock execution time.

All experiments were conducted on a 64-bit system with an Intel Core~i7 CPU (2.8~GHz) and 16~GB of RAM. The algorithms were implemented in Python~3.11.9 and executed sequentially.
